% ============================================================
% Conclusions — Paper A (JNCA)
% ============================================================
\section{Conclusions}
\label{sec:conclusion}

Modern phishing detection works in aggregate, yet a residual fraction---spear phishing and BEC---remains the
most costly and largely evades content filters, and lateral phishing from a compromised account is its hardest
case, because every static relation the sender exhibits is genuine. The gap is temporal. By modeling the
organization's communication graph with its rhythm, the violation becomes measurable in the dynamics.

We answer the two research questions in turn.
\begin{description}[noitemsep,topsep=2pt]
  \item[RQ1 (signal, T1):] A compromised account consistent with every static relation is detected from the
  receiving-side wave of its cascade. Unlike volume-based detection, the temporal detector
  stays above the volumetric one at every stealth level, though both degrade ($0.57$ vs $0.50$ at $K{=}1$;
  Section~\ref{sec:res-stealth}); the
  advantage vanishes under a time-shuffle control; the hierarchy replicates on Enron in both regimes and the
  regime contrast on two SNAP graphs (Section~\ref{sec:res-realgraphs}).
  \item[RQ2 (limits, T2):] The signal needs fan-out: a single-recipient attack is invisible to every
  detector (Section~\ref{sec:res-baselines}). The advantage degrades gracefully with scale and shrinks under
  label noise, and at a realistic base rate the detectors act as triage instruments, with effectiveness defined
  operationally through low-FPR sensitivity under ablation controls and explicit prevalence arithmetic
  (Section~\ref{sec:res-rigor}).
\end{description}

As its contributions, the paper offers a formalization of lateral phishing as a temporal violation of the
organizational communication graph, with a cascade threat model that includes the stealthy attacker; the
receiving-side cascade signal as a detection invariant, with a hierarchy of detectors, a time-shuffle control, and
validation on three real communication graphs; and a leak-aware evaluation methodology with an operational
definition of detection effectiveness for the targeted tail.

Among the limitations, the maliciousness is synthetic, so we validate the mechanism rather than the detection
of real intrusions; the cascade parameters are ours; and both scale and the operational threshold remain open.
Section~\ref{sec:discussion} gives the details.

% ZMIANA: nowe liczby po przepisaniu licznikow wolumetrycznych na przyczynowosc zdarzeniowa:
% zwykla kaskada 0.88, atak rozlozony w czasie 0.57-0.76, detektor wolumetryczny 0.50-0.63
% (blisko losowosci w reżimie skrytym; w blast widzi juz burst, ale to przypadek, ktorego
% atakujacy moze trywialnie uniknac).
The results turn lateral phishing into a quantity that can be tracked rather than merely flagged. Under the
default cascade the receiving-side signal reaches AUC $0.88$; under a stealthy attacker it degrades to
$0.57$--$0.76$ but stays above volume, which remains near chance in the stealthy regime ($0.50$--$0.64$) and
only recovers under the naive blast;
and the same regime transition appears in direction on Enron (now in both regimes), \texttt{email-Eu-core}, and
\texttt{CollegeMsg}. The defense this points toward
asks not whether a message looks malicious, but whether the organization's communication graph is behaving
the way it usually does. Turning that question into an operational filter in the SOC pipeline, at the
precision real deployment demands, remains an open problem.
